\topsection[sec:soa]{State of the art and critical analysis}

This section is the summary of phase~2. The full synthesis is \repofile{docs/02_state_of_the_art.md} and the
complete comparative table is \repofile{docs/03_comparative_table.md}. Every reference was opened and checked
in November 2026.

\subsection{Why an AI supply chain is not a software supply chain}
In a software supply chain, readable source code goes through a build that can be made deterministic. Anyone can
rebuild the binary and compare it with the one that was shipped, so reproducible builds and build provenance
give strong guarantees about what a binary contains. A model is produced differently. Its behaviour comes from
data, code, hyper-parameters and randomness. The result, a file of weights, cannot be read like source code,
training is rarely reproducible bit for bit, and the only way to learn what a model does is to test it. Malice
hides in a statistical property of the weights, not in lines that a reviewer can read, and a backdoored model
passes ordinary accuracy tests. Its upstream dependencies are not only libraries but datasets, pretrained
models and fine-tuning services.

Two consequences shape the project. Provenance must cover the whole lineage of a model, not one build. And even
complete provenance proves who produced which bytes; it does not prove that the bytes behave well.

\subsection{Attacks on AI supply chains}
\textbf{Direct poisoning of training and fine-tuning data.} BadNets \cite{badnets} showed that a few training
images stamped with a small pixel pattern and relabelled are enough to plant a backdoor: the model keeps its
normal accuracy and outputs the attacker's class whenever the pattern is present. Poisoning has since become
cheap at scale. Carlini et al. poisoned web-scale datasets for tens of dollars by buying expired domains that
dataset indexes still referenced \cite{carlini2024}. Souly et al. backdoored language models of 600 million to
13 billion parameters with a near-constant number of about 250 poisoned documents, which means that a larger
training set does not dilute the attack \cite{souly2025}. In every case the poison is present before anyone
signs the data.

\textbf{Backdoored and malicious pretrained models.} A model from a public hub can carry a backdoor in its
weights, or attack the machine that loads it. Python's pickle format, used by \code{torch.save}, executes code
during deserialisation. JFrog found about one hundred such models on Hugging Face in February 2024, one of them
opening a reverse shell \cite{jfrog2024}. In February 2025 ReversingLabs described models whose deliberately
broken pickle stream ran its payload before the security scanner failed \cite{rl2025}. A signature on such a
file only proves who published it.

\textbf{Environment poisoning.} The code and the infrastructure around training are part of the supply chain.
Between 25 and 30 December 2022 the \code{torchtriton} dependency confusion exfiltrated SSH keys and files from
machines that installed the PyTorch nightly builds \cite{pytorch2022}. A compromised training host can alter
weights during training and, if it holds a signing key, sign the result.

\textbf{Regulation and risk frameworks.} Besides Article 15(5) of the EU AI Act \cite{aiact2024}, the NIST AI
Risk Management Framework \cite{nistairmf}, the MITRE ATLAS technique AML.T0010 on AI supply chain compromise
\cite{atlas} and the OWASP Machine Learning Security Top~10 (ML06:2023) \cite{owaspml} describe the same risk.
None prescribes a mechanism. They ask organisations to be able to show where a model came from and that it was
not manipulated.

\subsection{Cryptographic signing and attestation}
A cryptographic hash such as SHA-256 gives a short fingerprint of a file: any change to the file changes the
fingerprint. A digital signature over that fingerprint proves that the holder of a private key approved exactly
these bytes. Ed25519 \cite{rfc8032} gives small keys, small signatures and fast verification. Two standards
define how such evidence is packaged: the DSSE envelope, which signs a typed payload \cite{dsse}, and the in-toto
attestation framework, whose statements bind a subject (a file digest) to a predicate such as build provenance
\cite{intotoattest}.

Key management is the hard part of signing. Sigstore removes long-lived keys: Fulcio issues short-lived
certificates bound to an OpenID Connect identity, Rekor records every signature in a public transparency log,
and \code{cosign} signs and verifies \cite{sigstore2022}. OpenSSF Model Signing 1.0, released in April 2025,
applies this to machine-learning models and signs a model directory file by file, with Sigstore or a private
PKI \cite{oms2025}. These tools prove the integrity and the origin of one artifact. They do not link a fine-tuned
model to the dataset and the base model it was trained from, they do not say which identity may sign which
stage, and they say nothing about what a model learned.

\subsection{Provenance frameworks and tamper-evident logs}
SLSA 1.1 defines levels of build provenance, from documented builds to provenance that a hardened builder
produces and that cannot be forged \cite{slsa11}. Applied to machine learning, training plays the role of the
build and the dataset and the base model are build inputs, but SLSA describes one build step, not a chain of
training and fine-tuning steps owned by different parties. in-toto is the closest conceptual match to this
project: a project owner signs a layout that lists the steps of a supply chain, who may perform each step and
which materials and products link the steps \cite{intoto2019}. Its layouts were designed for software builds and
do not model dataset content or revocation along an AI lineage.

The Update Framework showed how to survive a key compromise with separate roles, thresholds and explicit
revocation \cite{tuf2010}. Certificate Transparency showed how a Merkle tree makes a log append-only and
auditable \cite{rfc6962}. A Merkle tree hashes items in pairs up to a single root, so one root commits to every
item, and a changed item can be located. We use it to commit a dataset row by row.

\textbf{Blockchain anchoring}, which the brief lists as optional, gives a public, tamper-evident timeline at
the cost of operating or paying for a ledger and of trusting its consensus. For one organisation, a signed,
hash-chained log with a pinned head gives the same append-only property with far less infrastructure. A public
anchor becomes useful when several organisations must agree on one history. We therefore built the signed log
and left anchoring as future work.

\subsection{Documentation, content-level defences and safe formats}
Model cards \cite{modelcards} and datasheets for datasets \cite{datasheets} document intended use, training data
and evaluation. Bills of materials for machine learning list the components of an AI system: CycloneDX added
this in version 1.5 \cite{cyclonedx15} and SPDX 3.0 introduced AI and Dataset profiles \cite{spdx3}. These
documents bring transparency, not integrity. They are unsigned by default.

Content-level defences look at what the data or the model contains: data validation and label-distribution
tests on the data side; activation clustering \cite{activation}, spectral signatures \cite{spectral}, Neural
Cleanse \cite{neuralcleanse} and STRIP \cite{strip} on the model side. All of them give evidence, not proof.
Adaptive attackers choose triggers that a given detector does not model, and every detector has false positives
and a cost. Finally, the pickle attacks above are prevented by formats that cannot execute code. safetensors
stores a JSON header and raw tensors, and was audited by Trail of Bits in 2023 \cite{safetensors}.

\subsection{Open-source, industrial and academic approaches compared}
Table~\ref{tab:compare} compares the approaches capability by capability. No row reaches ``yes'' in the column
for poisoned content. Integrity mechanisms and content defences answer different questions, and a credible
control combines both and says which one produced each conclusion.

{\tablefont\setlength{\tabcolsep}{4pt}
\begin{xltabular}{\textwidth}{P{3.15cm}P{2.4cm}P{2.0cm}P{2.5cm}P{1.95cm}L}
\caption{Comparison of approaches (capabilities as documented by each project, November 2026; ``no'' means that
it is not a goal of the tool, not a defect).}\label{tab:compare}\\
\toprule
\headrow\hd{Approach} & \hd{Detects tampering after signing} & \hd{Proves origin} & \hd{Links artifacts across stages} & \hd{Detects poisoned content} & \hd{Needs} \\
\midrule
\endfirsthead
\toprule
\headrow\hd{Approach} & \hd{Detects tampering after signing} & \hd{Proves origin} & \hd{Links artifacts across stages} & \hd{Detects poisoned content} & \hd{Needs} \\
\midrule
\endhead
\bottomrule
\endfoot
\multicolumn{6}{l}{\textit{Building blocks}} \\*
Plain hashes & yes, if the hash comes from a trusted channel & no & no & no & nothing \\
Digital signatures & yes & yes (key holder) & no, unless parents are signed into the payload & no & key management \\
\midrule
\multicolumn{6}{l}{\textit{Open-source signing and provenance}} \\*
Sigstore, cosign & yes & yes (OIDC identity) & no (per artifact) & no & Fulcio, Rekor, identity provider \\
OpenSSF Model Signing & yes (model directory) & yes & no & no & Sigstore or a PKI \\
in-toto layouts & yes & yes & yes (layout of steps) & no & nothing \\
SLSA provenance & yes & yes (builder) & partial (build inputs) & no & hardened builder \\
TUF, transparency logs & yes (history, rollback) & yes & no & no & log service \\
\midrule
\multicolumn{6}{l}{\textit{Industrial practice and standards}} \\*
Model hubs: safe formats and file scanning & no & no & no & code on load only & hub \\
Model registries (MLflow, DVC) & no (metadata) & no (platform login) & lineage metadata & no & server \\
Model cards, ML-BOM (CycloneDX, SPDX) & only if signed & only if signed & yes (inventory) & no & nothing \\
Blockchain anchoring & yes (history) & partial & yes & no & ledger network \\
\midrule
\multicolumn{6}{l}{\textit{Academic defences}} \\*
Backdoor detectors, data validation & no & no & no & partial (heuristic, evadable) & nothing \\
\midrule
\textbf{\sys} & \textbf{yes}, to the row for datasets & \textbf{yes}, keys bound to roles & \textbf{yes}, parent digests inside each signed payload & partial, reported as \st{WARNING} only & keys in files; no outside service \\
\end{xltabular}}

\subsection{Limits of the state of the art and the approach we chose}
\begin{enumerate}
\item \textbf{No end-to-end lineage.} Existing tools sign one file or one build. Nothing links dataset, base
  model, fine-tuned model and deployment into one verifiable graph, decides which role may sign which stage,
  revokes along the links, or gates deployment on the whole chain.
\item \textbf{Unsigned metadata.} Lineage kept as ordinary database state can be edited by whoever
  administers the database.
\item \textbf{Signed is read as safe.} A valid signature is often presented to users as a safety signal. We
  found no work that measures how much protection it gives against poison that was present before signing.
\item \textbf{Content checks are partial.} They can be evaded and must be reported as separate, weaker
  evidence.
\end{enumerate}

\sys\ keeps the proven building blocks and adds what is missing for an AI lineage. Table~\ref{tab:choices}
gives each choice with its justification.

{\tablefont
\begin{xltabular}{\textwidth}{P{5.6cm}L}
\caption{The approach we selected and why.}\label{tab:choices}\\
\toprule
\headrow\hd{Choice} & \hd{Justification} \\
\midrule
\endfirsthead
\toprule
\headrow\hd{Choice} & \hd{Justification} \\
\midrule
\endhead
\bottomrule
\endfoot
Ed25519-signed JSON manifests with a canonical serialisation & Small, fast, deterministic signatures; one manifest per artifact; exportable as in-toto statements in DSSE envelopes \\
Each manifest embeds the manifest digests of its parents & Separate signatures become one hash-linked graph: rewriting an ancestor breaks every descendant, even with a legitimate key \\
Row-level Merkle root for datasets (RFC 6962 construction) & Detects and locates a single changed row \\
Signers bound to roles, and a stage policy & The in-toto layout idea applied to AI stages: the training key cannot sign datasets, and a fine-tuned model needs a dataset and a base model \\
Signed, hash-chained, pinned trust log for keys and revocation & Append-only history and revocation without operating a blockchain; the database is only an untrusted index \\
Deployment gateway that verifies the whole chain and loads the verified bytes; safetensors only & Closes the gap between the check and the use, and the code-execution path of pickle \\
Content checks reported separately, as \st{WARNING} & Keeps cryptographic and content evidence apart, which is the point the project measures \\
Custom signing at the core instead of Sigstore & Works offline, with the role policy and the parent links above; an optional Sigstore keyless job exists in the CI workflow (off by default) \\
\end{xltabular}}
